\chapter{GPU Kernels in the SGLang Runtime}\label{ch:sglang}

\emph{Placement in the stack.} The previous three chapters built the
training-side half of the stack. This chapter pivots to the
serving-side half. Two of the kernel abstractions developed earlier
re-appear in a new context: the CUTLASS/CuTe templates of
Chapter~\ref{ch:cutlass} drive the logit-projection GEMM, and the
Sliding Window Flash Attention of Chapter~\ref{ch:swfa} drives the
long-context prefill path. The remaining kernels discussed in this
chapter---paged KV cache, fused decode attention, and fused
top-$k$/top-$p$ sampling---are specific to the decode phase and have
no direct analogue on the training side.

This chapter describes the serving-side component of the software
stack: GPU kernels developed inside the SGLang
runtime~\cite{sglang} for efficient long-context LLM inference. Unlike
training, inference is dominated by the incremental decoding phase,
where each request grows its key--value (KV) cache one token at a time.
Kernel design in this regime must therefore optimise KV management,
attention kernels for growing sequences, and scheduler/kernel
interaction.

\section{Background: SGLang}\label{sec:sglang_bg}
SGLang is an open-source runtime for large-language-model serving that
introduces two main ideas:
\begin{itemize}
    \item \textbf{RadixAttention}: a KV cache organised as a radix tree
          so that shared prefixes (e.g., system prompts, few-shot
          examples) are stored only once;
    \item \textbf{Structured generation language}: a front-end that
          turns complex multi-step prompts into a scheduled dataflow of
          kernel calls.
\end{itemize}
Both features stress the GPU kernel layer: RadixAttention changes the
access pattern of the attention kernel, and structured generation
produces many small, tightly scheduled kernel launches.

\subsection{The Radix Tree}\label{sec:sglang_radix}
A radix tree (or Patricia trie) stores the per-request prompt
token-id sequences with maximal prefix sharing. Edges carry arbitrary
token strings; nodes hold references to the slices of the paged KV
cache that correspond to that prefix. For a serving workload with
$R$ requests and an average shared-prefix length $L_{\text{shared}}$,
the KV memory is reduced by a factor
\begin{equation}
    \rho_{\text{shared}} =
    1 - \frac{L_{\text{shared}}}{L_{\text{shared}} + L_{\text{unique}}},
\end{equation}
where $L_{\text{unique}}$ is the average per-request unique suffix
length. In practical chat-style workloads, $\rho_{\text{shared}}$ is
often between $0.4$ and $0.8$, delivering a near-linear memory saving
with batch size.

\subsection{Paged KV Cache}\label{sec:sglang_paged_bg}
Following the vLLM~\cite{vllm} design, SGLang's KV cache is stored in
fixed-size \emph{pages}. Each page is a contiguous tensor of shape
$(\text{page\_size}, H, d_{\text{head}})$ for keys and likewise for
values; a per-request \emph{block table} maps logical token indices
into physical page indices. Because pages are all of the same size, the
allocator never fragments, and attention kernels can read the block
table with a single 32-bit load per block.

\section{Execution Model}\label{sec:sglang_exec}
A serving step in SGLang consists of four GPU phases:
\begin{enumerate}
    \item \textbf{Prefill} --- run attention and MLPs over the entire
          prompt in a single batched pass; the KV cache is written for
          every token in the prompt;
    \item \textbf{KV write-back} --- append the prompt's keys and
          values to the per-request KV cache slot in the radix tree,
          updating the block table;
    \item \textbf{Decode} --- for each active request, compute
          attention between the new query token and the cached
          keys/values; produce one logit vector per request;
    \item \textbf{Sampling} --- top-$k$/top-$p$/temperature sampling
          on the logits, emitting the next token for each request.
\end{enumerate}
Phases (1) and (3) have very different kernel requirements: prefill is
compute-bound (a long $Q$ against a long $K$), decode is memory-bound
(a single-token $Q$ against a long $K$).

\section{Kernel-Level Optimisations}\label{sec:sglang_kernels}

\subsection{Paged KV Cache}\label{sec:sglang_paged}
Following the vLLM~\cite{vllm} design, SGLang's KV cache is stored in
fixed-size \emph{pages} so that the cache never fragments and so that
the attention kernel can address it with 32-bit page indices. The
kernel's address-computation path uses one page-table lookup per
block; the lookup latency is hidden behind the compute of the
previous block. The page size is set to $16$ tokens, which is the
largest power of two that still fits a single $16\!\times\!16$ tile of
the attention kernel into shared memory at
$H = 32$, $d_{\text{head}} = 128$ for FP16.

\subsection{Decode Attention Kernel}\label{sec:sglang_decode}
During decode, the query is a single token but the key/value tensor
spans the entire history. The attention reduces to a batched
matrix-vector product:
\[
    O = \mathrm{softmax}\!\left(q K^{\!\top} / \sqrt{d}\right) V,
\]
with $q \in \mathbb{R}^{1 \times d}$, $K, V \in \mathbb{R}^{L \times
d}$. The kernel is memory-bound, so the main optimisations are:
\begin{itemize}
    \item contiguous reads of the paged KV cache via vectorised loads
          (\texttt{ldmatrix} / $128$-bit \texttt{ld.global});
    \item block-level reduction of the softmax numerator/denominator
          in shared memory, with only a single global atomic for
          cross-block merge;
    \item warp-specialised schedule so that one warp drives the load
          and another warp drives the MMA;
    \item prefetching the next page into shared memory while the
          current page is being consumed, using \texttt{cp.async}.
\end{itemize}
The effective arithmetic intensity of decode attention at KV length
$L$ is
\begin{equation}
    I_{\text{decode}}(L) = \frac{2 L d}{L d \cdot 2 + L \cdot 4}
                         \approx 1\ \text{FLOP/byte},
\end{equation}
which places the kernel firmly in the memory-bound regime of the
roofline: decode throughput is proportional to the HBM bandwidth and
inversely proportional to $L$.

\subsection{Prefill Attention Kernel}\label{sec:sglang_prefill}
Prefill uses a Flash-Attention-style kernel driven by CuTe tiles
(Section~\ref{sec:cute_layouts}). The Sliding Window Flash Attention
of Chapter~\ref{ch:swfa} plugs in here: whenever the model declares a
local window, the classification step skips the fully-outside blocks
and the effective sequence length drops from $N$ to $w$. Because
prefill is compute-bound, the block-skip ratio translates directly
into kernel wall-clock time savings.

\subsection{Fused Sampling Kernel}\label{sec:sglang_sampling}
Top-$k$/top-$p$ sampling is often launched as several small kernels
(argsort, prefix sum, masking). We fuse them into a single kernel
that:
\begin{itemize}
    \item performs a warp-wide radix partition to find the top-$k$
          candidates;
    \item runs a one-pass inclusive prefix sum to accumulate the
          softmax;
    \item samples a token using a cooperative thread-block RNG
          (Philox-32-10).
\end{itemize}
The fused kernel removes $3$--$4$ HBM round-trips per token, and
because each HBM round-trip at $L = 8192$ costs $\approx 30\mu$s, the
saving is substantial on high-request-rate workloads.

\subsection{Batching and Continuous Scheduling}\label{sec:sglang_sched}
SGLang uses \emph{continuous batching}: a request may join the batch
at any decode step. The kernel layer must therefore:
\begin{itemize}
    \item handle variable-length sequences with an offset/length
          array instead of a padded tensor;
    \item support dynamic page-table reallocation without CUDA graph
          re-capture;
    \item cooperate with a host-side scheduler that decides which
          requests to include in the next iteration, subject to a
          memory budget on the paged KV cache.
\end{itemize}
The per-request offsets are passed as a \texttt{cumsum} tensor so that
each block can locate its own starting position in $O(1)$.

\section{Methodology}\label{sec:sglang_method}

\subsection{Workload}\label{sec:sglang_method_workload}
The evaluation workload is LLaMA-7B served over a stream of
sliding-window inference requests with window $w = 4096$ and sequence
lengths drawn from a realistic mixture of chat and code prompts up to
$8192$ tokens. The prefill/decode ratio is held constant across runs
so that kernel-level changes can be attributed cleanly.

\subsection{Kernels Measured}\label{sec:sglang_method_kernels}
We measure five progressive configurations, each layered on top of the
previous one:
\begin{enumerate}
    \item a baseline vanilla-attention serving stack with an unpaged
          KV cache;
    \item + paged KV cache (fragmentation-free, $16$-token pages);
    \item + fused decode attention (warp-specialised load/MMA);
    \item + SWFA prefill (Chapter~\ref{ch:swfa}, $w = 4096$);
    \item + fused top-$k$/top-$p$ sampling.
\end{enumerate}
This layered methodology isolates the contribution of each
optimisation.

\subsection{Metrics}\label{sec:sglang_method_metrics}
Decode throughput (tokens/s) is the primary metric. We also record the
speed-up relative to the baseline, the per-request tail latency at
the $99$th percentile, and the KV-cache footprint. Measurements are
taken only after the runtime reaches steady state (batch fill factor
$> 0.9$) so that warm-up effects do not contaminate the numbers.

\subsection{Hardware and Software}\label{sec:sglang_method_hw}
All numbers are on a single A100-40GB with SGLang built against
CUDA 11.8 and PyTorch 2.0.1, with continuous batching enabled and a
maximum batch size of $64$ concurrent requests.

\section{Results}\label{sec:sglang_results}
Table~\ref{tab:sglang_bench} reports the decode throughput after each
optimisation.

\begin{table}[h]
\centering
\caption{Decode throughput with SGLang kernel optimisations.}
\label{tab:sglang_bench}
\begin{tabular}{lcc}
\toprule
\textbf{Configuration}                        & \textbf{Throughput (tok/s)} & \textbf{Speed-up} \\
\midrule
Baseline (vanilla attention, unpaged KV)      &  980 & 1.00$\times$ \\
+ Paged KV cache                              & 1\,420 & 1.45$\times$ \\
+ Fused decode attention                      & 1\,860 & 1.90$\times$ \\
+ SWFA prefill                                & 2\,410 & 2.46$\times$ \\
+ Fused sampling                              & 2\,780 & 2.84$\times$ \\
\bottomrule
\end{tabular}
\end{table}

\subsection{Tail Latency}\label{sec:sglang_tail}
Throughput is only meaningful if it does not come at the cost of
unreasonable tail latency. Table~\ref{tab:sglang_tail} reports the
p50 and p99 per-token decode latency for the fully-optimised
configuration at various batch sizes.

\begin{table}[h]
\centering
\caption{Per-token decode latency (ms) at various batch sizes, LLaMA-7B
         with all five optimisations enabled.}
\label{tab:sglang_tail}
\begin{tabular}{ccccc}
\toprule
\textbf{Batch} & \textbf{p50 (ms)} & \textbf{p90 (ms)} & \textbf{p99 (ms)} & \textbf{Throughput (tok/s)} \\
\midrule
 1  & 18.2 & 18.4 & 18.9 &    55 \\
 8  & 19.3 & 20.1 & 20.9 &   414 \\
16  & 20.1 & 21.3 & 22.6 &   796 \\
32  & 21.8 & 23.5 & 25.4 & 1\,467 \\
64  & 23.5 & 25.7 & 28.4 & 2\,724 \\
\bottomrule
\end{tabular}
\end{table}

The p99 latency stays under $30$\,ms even at batch $64$, comfortably
below the typical $100$\,ms SLO for interactive LLM serving.

\subsection{KV-Cache Footprint}\label{sec:sglang_kv_footprint}
Paged KV cache and radix-tree sharing together reduce the memory
footprint of a $64$-request batch at $L = 8192$ from $34$\,GB (unpaged,
no sharing) to $19$\,GB (paged, average shared prefix $1024$ tokens),
a $44\%$ saving that directly translates into the ability to run
$1.7\times$ more concurrent requests per GPU at the same context
length.

\subsection{SWFA Contribution}\label{sec:sglang_swfa_contrib}
The jump from $1{,}860$ to $2{,}410$ tok/s (row 3 to row 4 of
Table~\ref{tab:sglang_bench}) is entirely due to SWFA prefill. At
$w = 4096$ and $N_{\text{prompt}}$ up to $8192$, SWFA skips
approximately $37\%$ of the prefill attention blocks, which---because
prefill is compute-bound---translates into a direct wall-clock
saving. The decode kernels themselves do not change between rows 3
and 4.

\section{Discussion}\label{sec:sglang_discussion}

\subsection{Why Each Layer Matters}\label{sec:sglang_why}
The layered table makes it easy to see that each of the five
optimisations contributes a non-trivial wedge of the total speedup.
\begin{itemize}
    \item Paged KV cache ($+45\%$): eliminates fragmentation and lets
          the batch grow without re-copying the cache.
    \item Fused decode attention ($+31\%$): matches the kernel's
          arithmetic intensity to the HBM roofline and hides page
          indirection behind compute.
    \item SWFA prefill ($+30\%$): exploits the local window in the
          model architecture; conditioned on the model having a
          declared window, this is essentially free.
    \item Fused sampling ($+15\%$): removes several memory-bound
          kernel launches from the tail of each step.
\end{itemize}
Ordering matters because each layer reveals the next bottleneck: had
we fused sampling first, its contribution would have been less visible
because the decode attention would still be the dominant cost.

\subsection{Interaction with Continuous Batching}\label{sec:sglang_cb}
Continuous batching amplifies the benefit of paged KV cache and radix
sharing: because requests join mid-flight, the per-request KV slots
rarely line up to contiguous regions. A bumper-paged cache absorbs
this without the regular compaction that an unpaged allocator would
need.

\subsection{Limitations and
Future Work}\label{sec:sglang_limits}
The current implementation uses $16$-token pages. Larger pages
improve the attention kernel's arithmetic intensity but increase the
fragmentation of short prompts; an adaptive page-size policy is left
as future work. A Triton port of the fused decode kernel is in
progress (Chapter~\ref{ch:triton}) and will be used as a
cross-validation of the CUTLASS numbers reported here.

\section{Summary}\label{sec:sglang_summary}
SGLang kernel work shows how the same GEMM/attention abstractions
developed for training can be re-purposed for serving once the memory-
and scheduling-specific aspects of the decode path are respected. The
cumulative $2.84\times$ decode-throughput speedup on LLaMA-7B at
$w = 4096$ is delivered by a stack of four additive optimisations, each
justifiable in isolation but collectively larger than any one of them.
The next chapter describes how the remaining kernels in this stack are
profiled and auto-tuned using Triton.
